\scanpage{50}
algébriques sur $X\times Y$, via les identifications \unsure{du|des} type (pour une variété $T$)
\[
\left\{\begin{aligned}
H^1(A_T,\mathrm{Q}_{\ell})&\xrightarrow{\sim}H^1(T,\mathrm{Q}_{\ell})\\
H^{\deleted{\unsure{2n-|2n-1}}1}(B_T,\mathrm{Q}_{\ell})&\xleftarrow{\sim}H^{2n-1}(T,\mathrm{Q}_{\ell}\deleted{\unsure{(}})(n-1)
\end{aligned}\right.
\]
De plus, l'hom correspondant entre groupes de coh. $l$-adiques relatifs à $X$ et à $Y$, est un isom ssi le $\mathrm{Q}$-hom \deleted{correspondant} envisagé $A\unsure{-|\longrightarrow}B$ est une \uline{isogénie}. Cela prouve en particulier que dans ce cas, il y a une classe de corr. $\mathrm{Q}$-algébriques sur $X\times Y$ définissant l'hom inverse sur la cohomologie.

Nous verrons que cette circonstance se produit en particulier si on a des classes de diviseurs amples $\unsure{\zeta|\xi}_1,\ldots,\zeta_{n-1}$ sur $X$ de dim $n$, en prenant l'hom de cup-produit par \deleted{$\unsure{\xi|\zeta}$} $\zeta_1\cdot\ldots\zeta_{n-1}:H^1(X,\mathrm{Q}_{\ell})-H^{2n-1}(X,\mathrm{Q}_{\ell})(n-1)$.
% archive sheet number: 49
